\documentclass[11pt,a4paper]{article}
\usepackage[utf8]{inputenc}
\usepackage[T1]{fontenc}
\usepackage{lmodern}
\usepackage[portuguese]{babel}
\usepackage[a4paper,margin=2.1cm]{geometry}
\usepackage{xcolor}
\usepackage{graphicx}
\usepackage{titlesec}
\usepackage{tcolorbox}
\usepackage{amsmath}
\usepackage{fancyhdr}
\usepackage{hyperref}
\tcbuselibrary{skins,breakable}

\definecolor{BgCream}{HTML}{FAF8F5}
\definecolor{TextEspresso}{HTML}{4A4238}
\definecolor{NudeTaupe}{HTML}{BFAFA6}
\definecolor{SoftBlush}{HTML}{D9C5B2}

\hypersetup{colorlinks=true, linkcolor=TextEspresso, urlcolor=TextEspresso}

\pagestyle{fancy}
\fancyhf{}
\renewcommand{\headrulewidth}{0pt}
\fancyfoot[C]{\color{NudeTaupe}\small\thepage}
\setlength{\parindent}{0pt}
\setlength{\parskip}{0.32cm}

\newtcolorbox{slidebox}{
  breakable, colback=BgCream, colframe=NudeTaupe,
  boxrule=0.6pt, arc=2pt, left=4pt, right=4pt, top=4pt, bottom=4pt
}

\newcommand{\sebentatitle}[3]{%
  \thispagestyle{empty}
  \begin{center}
  {\color{NudeTaupe}\rule{\textwidth}{0.5pt}}\\[0.5cm]
  {\LARGE\bfseries\color{TextEspresso} Sebenta #1}\\[0.35cm]
  {\large\color{TextEspresso} #2}\\[0.3cm]
  {\normalsize\color{NudeTaupe} Infraestruturas de Computação na Nuvem \quad|\quad #3}\\[0.4cm]
  {\color{NudeTaupe}\rule{\textwidth}{0.5pt}}
  \end{center}
  \vspace{0.4cm}
}

\newcommand{\pagina}[2]{%
  \begin{slidebox}
  \centering
  \includegraphics[width=0.95\textwidth]{#1}
  \end{slidebox}
  \vspace{0.4cm}
  #2
  \clearpage
}

\begin{document}

\sebentatitle{02}{Arquitetura de Datacenters e Modelos de Operação}{Aula 02}

\section*{Objeto e utilização deste documento}

O presente documento acompanha a apresentação da Aula 02. Cada página reproduz um slide e desenvolve, em texto corrido, a matéria nele condensada. A aula integra sete questões de escolha múltipla, cuja resolução é apresentada no slide imediatamente seguinte. Recomenda-se que a resposta seja formulada antes da leitura da resolução.

\clearpage

\pagina{slides/slide-01.png}{
A Aula 02 desloca a análise do plano conceptual estabelecido na Aula 01 para o objeto físico que sustenta a computação na nuvem: a instalação onde os sistemas são alojados. A exposição desenvolve-se em três partes de complexidade crescente. A primeira estabelece a definição de datacenter e os critérios normativos de classificação da sua redundância. A segunda trata da rede interna e da quantificação da disponibilidade. A terceira aborda os modelos de operação, a distribuição geográfica da infraestrutura e a eficiência energética. Sete questões de escolha múltipla distribuem-se pela exposição, com resolução fundamentada no slide seguinte.
}

\pagina{slides/slide-02.png}{
A primeira parte estabelece os fundamentos. Define o objeto de estudo, identifica os subsistemas que o compõem e introduz o vocabulário normalizado de redundância, do qual decorre a classificação em Tiers. Não pressupõe conhecimento prévio sobre infraestrutura física.
}

\pagina{slides/slide-03.png}{
A segunda parte trata da rede interna da instalação. Compara as duas topologias de referência, caracteriza os padrões de tráfego que justificam a preferência por uma delas, e introduz a razão de sobre-subscrição enquanto instrumento de dimensionamento das ligações ascendentes.
}

\pagina{slides/slide-04.png}{
A terceira parte procede à quantificação da disponibilidade a partir dos tempos médios característicos do sistema, estabelece os modelos de composição em série e em paralelo, analisa a estrutura formal de um acordo de nível de serviço, e conclui com os modelos de operação, a organização geográfica da infraestrutura de cloud e as métricas de eficiência energética.
}

\pagina{slides/slide-05.png}{
A definição merece precisão. Um datacenter é uma instalação destinada a alojar sistemas de computação e a respetiva infraestrutura de suporte, dimensionada para operação contínua. A caracterização releva sobretudo pelo que acrescenta à noção corrente de sala de servidores: o datacenter compreende quatro subsistemas distintos, a distribuição de energia, o controlo térmico, a conectividade de rede e o espaço físico com os respetivos controlos de acesso. Estes subsistemas encontram-se em dependência mútua, de tal modo que a indisponibilidade de qualquer um compromete a operação do conjunto. Esta interdependência constitui o fundamento de toda a análise subsequente, designadamente da classificação em Tiers e do cálculo de disponibilidade.
}

\pagina{slides/slide-06.png}{
O diagrama representa a articulação dos quatro subsistemas numa instalação elementar. A alimentação elétrica é estabilizada por UPS e grupo gerador antes de alcançar os bastidores; o controlo térmico, assegurado por unidades CRAC ou CRAH, mantém a sala dentro dos limites de operação do equipamento; a conectividade de rede estabelece a ligação ao exterior. Cada ligação representada no diagrama constitui uma dependência, e a sua interrupção afeta a totalidade dos sistemas a jusante.
}

\pagina{slides/slide-07.png}{
Introduz-se aqui um conceito que estrutura toda a análise de arquitetura: o domínio de falha, definido como o conjunto de componentes afetados pela falha de um elemento comum de que todos dependem. Um quadro de distribuição elétrica, uma unidade de arrefecimento, um comutador de topo de bastidor ou um troço de fibra óptica delimitam, cada um, o seu domínio. Quando esse domínio coincide com a totalidade do sistema, o componente constitui um ponto único de falha. Importa notar o deslocamento de perspetiva que este conceito impõe: a análise de arquitetura não se ocupa de determinar se um componente pode falhar, questão cuja resposta se assume afirmativa, mas de determinar a extensão e a via de propagação dessa falha.
}

\pagina{slides/slide-08.png}{
A primeira questão examina a propagação de uma falha elétrica numa instalação desprovida de qualquer redundância. Recomenda-se a formulação de uma resposta, e eventual discussão, antes da consulta do slide seguinte.
}

\pagina{slides/slide-09.png}{
A fundamentação da resposta consta do slide. Importa reter a articulação com a matéria subsequente: a configuração descrita, na qual a via única de alimentação constitui um ponto único de falha de domínio total, corresponde exatamente ao nível mais baixo da classificação normativa apresentada em seguida. Os níveis superiores definem-se precisamente pelas vias alternativas que esta configuração não possui.
}

\pagina{slides/slide-10.png}{
Assumida a inevitabilidade da falha, o critério de projeto relevante passa a ser o número de vias alternativas disponíveis antes da interrupção do serviço. A designação normalizada destes níveis recorre a três expressões que reaparecem ao longo do semestre. O nível $N$ corresponde à capacidade estritamente necessária para sustentar a carga nominal, sem qualquer margem. O nível $N+1$ acresce uma unidade de reserva ao conjunto necessário, tolerando a indisponibilidade de um componente, mas mantendo via única de distribuição. O nível $2N$ procede à duplicação integral do subsistema, incluindo as vias de distribuição, e tolera a perda de metade da instalação sem degradação do serviço. Encontra-se ainda com frequência a variante $2(N+1)$, adotada quando se exige manutenção concorrente sobre uma via já redundante.
}

\pagina{slides/slide-11.png}{
O Uptime Institute mantém, desde a década de noventa, o sistema de classificação de utilização mais generalizada na caracterização de datacenters. O âmbito da classificação requer precisão: incide sobre a topologia da infraestrutura, isto é, sobre a organização dos subsistemas e o número de vias independentes até cada equipamento, e não sobre o desempenho, a dimensão ou a qualidade do equipamento instalado. O Tier I designa infraestrutura de capacidade básica, desprovida de redundância, na qual qualquer intervenção de manutenção ou avaria implica interrupção do serviço. O Tier II introduz componentes de capacidade redundantes, em configuração $N+1$, mantendo contudo uma via única de distribuição que permanece como ponto único de falha.
}

\pagina{slides/slide-12.png}{
A transição qualitativa verifica-se no Tier III, que dispõe de múltiplas vias de distribuição independentes, das quais uma se encontra ativa. Esta configuração admite manutenção concorrente, isto é, a intervenção sobre qualquer componente sem interrupção do serviço, o que os dois níveis anteriores não permitem. O Tier IV exige redundância $2N$ em todos os componentes críticos, acrescida de compartimentação física entre vias, de modo que uma falha única não produza efeito observável no serviço. Impõe-se um critério de aplicação frequentemente negligenciado: o nível atribuído a uma instalação corresponde ao do subsistema menos redundante, e não à média dos subsistemas. Este critério é objeto da questão seguinte.
}

\pagina{slides/slide-13.png}{
A segunda questão apresenta uma instalação com redundância assimétrica entre subsistemas e solicita a atribuição do nível correspondente. O enunciado contém deliberadamente elementos que sugerem uma classificação superior à correta.
}

\pagina{slides/slide-14.png}{
A fundamentação consta do slide anterior. O princípio subjacente admite formulação geral: a redundância de um sistema não corresponde à média das redundâncias dos seus constituintes, mas ao mínimo. Um único subsistema sem via alternativa determina o limite do conjunto, independentemente do grau de duplicação dos restantes. O mesmo princípio será aplicado, em contexto distinto, na análise de disponibilidade de serviços distribuídos.
}

\pagina{slides/slide-15.png}{
A classificação do Uptime Institute coexiste com outras normas de âmbito distinto. A ANSI/TIA-942, de origem norte-americana e enquadramento nas telecomunicações, define os níveis Rated-1 a Rated-4 e abrange adicionalmente matérias que os Tiers não detalham, designadamente cablagem estruturada, disposição do espaço e proteção contra incêndio. Na Europa é corrente a referência à série EN 50600, organizada em classes de disponibilidade de 1 a 4, que integra requisitos de eficiência energética e de segurança física numa abordagem de âmbito mais alargado. As escalas são análogas em princípio, distinguindo todas níveis crescentes de redundância, mas não estabelecem equivalência formal entre si. Em consequência, a caracterização de uma instalação deve explicitar sempre a norma aplicada.
}

\pagina{slides/slide-16.png}{
A topologia hierárquica de três camadas constituiu durante décadas a organização de referência para a rede interna de um datacenter, encontrando-se ainda em instalações de construção anterior. As três camadas desempenham funções distintas: a camada core assegura a ligação ao exterior e o encaminhamento entre blocos, a camada de aggregation concentra o tráfego e delimita domínios de difusão, e a camada de access estabelece a ligação física aos servidores. A limitação estrutural manifesta-se na comunicação entre servidores alojados em bastidores distintos, que exige a travessia ascendente e descendente da hierarquia, com acréscimo de latência proporcional ao número de saltos. Daqui decorre que a latência entre dois servidores depende da respetiva localização física, o que compromete a previsibilidade do desempenho.
}

\pagina{slides/slide-17.png}{
A topologia leaf-spine deriva das redes de Clos, formalizadas no contexto da comutação telefónica na década de cinquenta. A sua caracterização é simples: cada comutador leaf estabelece ligação a todos os comutadores spine, não existindo adjacências diretas entre leafs nem entre spines. Em consequência, o percurso entre quaisquer dois servidores ligados a leafs distintos compreende exatamente dois saltos, independentemente da localização física. A multiplicidade de percursos equivalentes permite ainda a distribuição de carga por ECMP, com agregação da largura de banda disponível entre os vários spines. A capacidade agregada escala pela adição de spines, sendo o número máximo de leafs determinado pela densidade de portas de cada spine.
}

\pagina{slides/slide-18.png}{
O diagrama apresenta as duas topologias em confronto. Na organização hierárquica, o percurso entre bastidores distintos exige a subida à camada de aggregation, por vezes à camada core, seguida da descida correspondente. Na organização leaf-spine, qualquer par de leafs dista exatamente dois saltos. A comparação visual explicita, sem recurso a formulação analítica, o fundamento da substituição progressiva da primeira pela segunda em instalações destinadas a suportar aplicações distribuídas.
}

\pagina{slides/slide-19.png}{
A distinção entre os dois padrões de tráfego refere-se à direção do fluxo relativamente à fronteira da instalação. O tráfego norte-sul atravessa essa fronteira, entre clientes externos e os sistemas alojados. O tráfego este-oeste circula exclusivamente no interior, resultando de invocações entre componentes de uma mesma aplicação. A decomposição de aplicações monolíticas em serviços distribuídos inverteu a proporção histórica entre os dois padrões: um único pedido externo pode originar dezenas de invocações internas antes de produzir resposta. É esta inversão que justifica a preferência pela topologia leaf-spine, cuja propriedade determinante consiste em uniformizar a distância entre quaisquer dois pontos internos da rede, e não apenas entre os sistemas internos e a saída para o exterior.
}

\pagina{slides/slide-20.png}{
A razão de sobre-subscrição quantifica a contenção admitida no dimensionamento de um comutador, definindo-se como o quociente entre a largura de banda agregada dos portos descendentes e a dos portos ascendentes. O exemplo apresentado no slide, um leaf com 48 portos de 10\,Gbit/s para servidores e 4 portos de 40\,Gbit/s para spines, conduz a uma razão de 3:1, o que significa que a capacidade de saída corresponde a um terço da capacidade de entrada. Uma razão unitária designa-se por non-blocking e assegura que qualquer padrão de tráfego é suportado sem contenção, ao custo de capacidade de uplink adicional. O valor admissível decorre do padrão de tráfego previsto: cargas com tráfego este-oeste intenso exigem razões próximas da unidade, ao passo que cargas dominadas por tráfego norte-sul toleram razões superiores.
}

\pagina{slides/slide-21.png}{
O caso proposto descreve uma plataforma de distribuição de vídeo decomposta em cerca de duzentos serviços, com volume de invocações internas substancialmente superior ao tráfego externo e sem padrão estável de comunicação entre pares. Recomenda-se a discussão da topologia adequada, e das consequências da alternativa, antes da consulta da questão seguinte, que retoma o mesmo cenário.
}

\pagina{slides/slide-22.png}{
A questão formaliza a decisão de projeto suscitada pelo caso anterior. A resposta deve fundamentar-se nas propriedades das duas topologias e no padrão de tráfego caracterizado no enunciado.
}

\pagina{slides/slide-23.png}{
A fundamentação consta do slide anterior. O elemento determinante reside na ausência de determinismo do padrão de comunicação: não sendo possível antecipar que pares de serviços comunicam, nem com que intensidade, a otimização da colocação física torna-se inviável, e a uniformização da latência assume-se como o único critério robusto. Acresce que o encadeamento de dezenas de invocações num único pedido externo amplifica qualquer variabilidade de latência individual no tempo de resposta agregado.
}

\pagina{slides/slide-24.png}{
A disponibilidade define-se como a fração do tempo em que o sistema presta serviço conforme a especificação, dada pelo quociente entre o tempo operacional e o tempo total. Em regime estacionário, a grandeza exprime-se em função dos tempos médios característicos do sistema, através da relação entre o MTBF, tempo médio entre falhas consecutivas, e o MTTR, tempo médio de restabelecimento do serviço. Esta formulação é preferível à definição empírica por permitir o cálculo prospetivo da disponibilidade a partir de parâmetros conhecidos do equipamento, sem necessidade de observação prolongada. Regista-se ainda a interpretação probabilística da grandeza: a disponibilidade corresponde à probabilidade de o sistema se encontrar operacional num instante arbitrário, propriedade que fundamenta a composição de componentes apresentada adiante.
}

\pagina{slides/slide-25.png}{
A expressão da disponibilidade em função do MTBF e do MTTR admite duas estratégias distintas de melhoria, correspondentes aos dois termos do quociente. O aumento do MTBF atua sobre a frequência das falhas, por seleção de componentes, redundância e controlo das condições de operação. A redução do MTTR atua sobre a duração da interrupção, por deteção automática, diagnóstico assistido, manutenção de peças de substituição em reserva e ensaio prévio dos procedimentos de recuperação. Em sistemas distribuídos, a segunda estratégia apresenta com frequência melhor relação custo-benefício, por não exigir substituição de hardware nem alteração da arquitetura física. A questão 5 quantifica precisamente este efeito.
}

\pagina{slides/slide-26.png}{
A tabela estabelece a correspondência entre as classes de disponibilidade correntes e o tempo de indisponibilidade que cada uma admite, em base anual e mensal. A leitura da tabela evidencia a natureza multiplicativa da progressão: cada ordem de grandeza adicional reduz por um fator de dez a indisponibilidade admissível. A passagem de 99,9\% a 99,99\% reduz o orçamento anual de aproximadamente oito horas e quarenta e seis minutos para cinquenta e dois minutos. A consequência operacional é determinante: a partir de 99,99\%, o orçamento anual não comporta intervenção manual, uma vez que a deteção, o diagnóstico e a recuperação teriam de concluir-se em poucos minutos. A partir desse limiar, a automatização dos procedimentos de recuperação deixa de constituir uma otimização para se tornar um requisito.
}

\pagina{slides/slide-27.png}{
Componentes dizem-se em série quando a totalidade é necessária ao funcionamento do sistema. Admitindo independência estatística das falhas, a disponibilidade do conjunto obtém-se pelo produto das disponibilidades individuais. Da propriedade $A_i \le 1$ decorre que o produto não excede o menor dos fatores, o que se traduz numa conclusão de alcance geral: a composição em série não excede a disponibilidade do componente menos disponível. Em consequência, o acréscimo de dependências em série degrada necessariamente a disponibilidade do conjunto, ainda que cada dependência considerada isoladamente apresente elevada fiabilidade. Esta observação justifica a atenção prestada, em arquitetura de sistemas, à redução do número de dependências obrigatórias no percurso de um pedido.
}

\pagina{slides/slide-28.png}{
Componentes dizem-se em paralelo quando é suficiente que um se mantenha operacional para que o sistema preste serviço. A indisponibilidade do conjunto exige, nesta configuração, a falha simultânea de todos os componentes, pelo que a disponibilidade se obtém por complemento do produto das indisponibilidades individuais. Verifica-se que o resultado excede a disponibilidade do componente mais disponível, o que estabelece a redundância como mecanismo de elevação da disponibilidade acima da de qualquer constituinte. A validade da expressão depende, contudo, da hipótese de independência das falhas. Componentes que partilham alimentação, via de rede, versão de software ou defeito de conceção apresentam correlação, e a disponibilidade efetiva situa-se abaixo do valor calculado. A verificação desta hipótese constitui parte integrante da análise.
}

\pagina{slides/slide-29.png}{
O diagrama representa as duas configurações. Na composição em série, os componentes dispõem-se em cadeia, e a interrupção em qualquer ponto do percurso compromete o conjunto. Na composição em paralelo, existem percursos alternativos entre origem e destino, e apenas a falha simultânea de todos interrompe o serviço. Os dois slides seguintes aplicam as expressões a valores concretos.
}

\pagina{slides/slide-30.png}{
O exemplo aplica a composição em série a um serviço dependente de um encaminhador e de uma instância de base de dados não replicada. O cálculo conduz a uma disponibilidade de aproximadamente 99,45\%, correspondente a cerca de quarenta e oito horas de indisponibilidade anual. Confirma-se a propriedade enunciada no slide 27: o resultado situa-se abaixo do menor dos dois fatores. A leitura arquitetural do resultado é a mais relevante: a base de dados não replicada, por apresentar a menor disponibilidade individual, determina o limite superior da disponibilidade do conjunto, e qualquer melhoria do encaminhador produziria efeito marginal enquanto essa dependência se mantivesse.
}

\pagina{slides/slide-31.png}{
O exemplo aplica a composição em paralelo a dois servidores aplicacionais equivalentes em configuração redundante. O procedimento de cálculo é inverso ao do caso anterior: determina-se a probabilidade de falha simultânea e subtrai-se à unidade. O resultado, 99,91\%, excede a disponibilidade individual de qualquer dos servidores, o que ilustra o efeito da redundância. O slide apresenta ainda o valor correspondente a três réplicas, 99,9973\%, que evidencia o retorno decrescente de cada réplica adicional: a segunda réplica produz um ganho muito superior ao da terceira, o que fundamenta a prática corrente de limitar a redundância a duas ou três instâncias e investir o remanescente na redução do MTTR.
}

\pagina{slides/slide-32.png}{
A quarta questão combina as duas configurações numa arquitetura mista, situação que corresponde à generalidade dos sistemas reais. Recomenda-se a resolução autónoma antes da consulta das opções apresentadas.
}

\pagina{slides/slide-33.png}{
A resolução consta do slide anterior. Importa assinalar a natureza do erro representado pela opção B, que resulta de aplicar a composição em série ao par redundante. O procedimento que o evita consiste em identificar a configuração de cada subsistema antes de qualquer cálculo, e proceder das composições redundantes mais internas para o exterior, reduzindo sucessivamente cada subsistema a um valor único de disponibilidade.
}

\pagina{slides/slide-34.png}{
A quinta questão quantifica o efeito da redução do MTTR sobre a disponibilidade, mantendo constante a frequência das falhas. O cálculo requer apenas a aplicação direta da expressão apresentada no slide 24, com atenção à uniformização das unidades de tempo.
}

\pagina{slides/slide-35.png}{
A resolução consta do slide anterior. O resultado fundamenta a orientação enunciada no slide 25: uma redução do MTTR por um fator de oito eleva a disponibilidade de 99,80\% para 99,98\%, o que corresponde a uma redução de uma ordem de grandeza na indisponibilidade anual, sem qualquer alteração de hardware nem da frequência das falhas. A opção B traduz o equívoco corrente de atribuir a disponibilidade exclusivamente à fiabilidade dos componentes, ignorando a contribuição do tempo de restabelecimento.
}

\pagina{slides/slide-36.png}{
O acordo de nível de serviço constitui o instrumento contratual que fixa o nível exigível ao fornecedor e as consequências do incumprimento. Os seus elementos constitutivos são cinco: o indicador medido e o respetivo método de medição, o valor de referência, o período de apuramento, as exclusões admitidas e o regime de compensação. A omissão de qualquer destes elementos compromete a exigibilidade do acordo, matéria retomada na questão 6. Regista-se ainda a articulação com a primeira parte da aula: o nível contratualizável é condicionado pela topologia da instalação, não sendo sustentável um compromisso de disponibilidade elevada a partir de uma configuração desprovida de redundância.
}

\pagina{slides/slide-37.png}{
Dois elementos determinam o valor efetivo de um acordo e merecem exame atento. As exclusões delimitam os períodos e as causas que não relevam para o apuramento, designadamente janelas de manutenção programada, causas imputáveis ao cliente e casos de força maior, pelo que um acordo com exclusões amplas pode apresentar um valor de referência elevado e cobertura reduzida. O regime de compensação assume tipicamente a forma de crédito de serviço proporcional ao desvio, e não de indemnização pelo prejuízo efetivamente incorrido. Daqui decorre uma limitação de alcance prático: o acordo transfere para o fornecedor uma fração reduzida do risco, pelo que não dispensa a construção de redundância do lado do cliente.
}

\pagina{slides/slide-38.png}{
A sexta questão examina a exigibilidade de compensação na ausência de indicador e de valor de referência contratualizados. A resposta decorre diretamente dos elementos constitutivos enunciados no slide 36.
}

\pagina{slides/slide-39.png}{
A fundamentação consta do slide anterior. O princípio admite formulação sintética: uma obrigação de melhor esforço não constitui obrigação de resultado, e na ausência de critério objetivo de incumprimento não existe facto gerador de compensação. Em termos práticos, a verificação dos cinco elementos constitutivos precede necessariamente a celebração do contrato.
}

\pagina{slides/slide-40.png}{
Os modelos de operação distinguem-se pela localização da fronteira de responsabilidade entre a organização e terceiros. No modelo on-premises, a organização detém e opera a totalidade da infraestrutura, assumindo o controlo integral da arquitetura e a responsabilidade pelos subsistemas de suporte analisados na primeira parte da aula. No modelo de colocation, a organização detém e opera o equipamento, contratando a um operador de instalação o espaço físico, a alimentação, o arrefecimento e a conectividade. No modelo de cloud pública, a infraestrutura é integralmente operada por terceiros, sendo o acesso mediado por interfaces de aprovisionamento programático. A fronteira adotada delimita simultaneamente o âmbito da responsabilidade operacional, o âmbito da responsabilidade de segurança e o grau de dependência técnica do fornecedor, aspeto retomado na questão 7.
}

\pagina{slides/slide-41.png}{
Os operadores de cloud estruturam a infraestrutura segundo uma hierarquia geográfica destinada a conter a propagação de falhas. As regiões constituem áreas geográficas distintas e mutuamente independentes quanto a alimentação e conectividade. Cada região subdivide-se em zonas de disponibilidade, instalações fisicamente separadas com subsistemas de suporte independentes, interligadas por ligações de latência reduzida que admitem replicação síncrona. A distribuição de réplicas por múltiplas zonas eleva o domínio de falha tolerado da instalação individual para a região; a distribuição por múltiplas regiões estende-o a eventos de escala geográfica. Regista-se uma restrição técnica relevante: a latência inter-regional inviabiliza, em geral, a replicação síncrona, impondo modelos de consistência mais fracos, matéria desenvolvida no bloco C5. A seleção da região encontra-se ainda condicionada por requisitos jurídicos de localização de dados.
}

\pagina{slides/slide-42.png}{
A sétima questão examina o alcance de uma garantia contratual de disponibilidade de zona relativamente a uma aplicação nela instalada. A resposta requer a aplicação do modelo de composição em série estabelecido no slide 27.
}

\pagina{slides/slide-43.png}{
A fundamentação consta do slide anterior e articula as duas metades da aula: a aplicação encontra-se em série com a zona, pelo que a disponibilidade observável não excede a da zona, independentemente da qualidade do software. A garantia do operador constitui, por conseguinte, um limite superior e não um valor assegurado. A elevação da disponibilidade acima desse limite exige composição em paralelo, mediante distribuição de réplicas por zonas independentes, o que transfere para o cliente a responsabilidade pela arquitetura de redundância.
}

\pagina{slides/slide-44.png}{
Duas tendências de implantação merecem registo, com critérios de projeto distintos. As instalações modulares consistem em unidades pré-fabricadas, integradas e validadas em fábrica antes da expedição, que reduzem o prazo de implantação de anos para semanas e permitem expansão incremental da capacidade pela adição de módulos. As instalações de periferia consistem em infraestrutura de dimensão reduzida, implantada em proximidade geográfica dos consumidores, com o objetivo de reduzir a latência de propagação. A distinção entre ambas reside no critério otimizado: as primeiras otimizam o prazo e o custo de expansão da capacidade, as segundas otimizam a latência, aceitando menor eficiência de escala.
}

\pagina{slides/slide-45.png}{
O consórcio The Green Grid normalizou um conjunto de indicadores de eficiência, dos quais o mais difundido é o PUE, definido pelo quociente entre a energia total consumida pela instalação e a energia consumida pelo equipamento de TI. Por construção, o indicador não assume valores inferiores à unidade. A sua interpretação é direta: um valor de 1,5 significa que, por cada unidade de energia consumida pelo equipamento de TI, meia unidade adicional é despendida em arrefecimento, conversão e iluminação. Instalações otimizadas operam abaixo de 1,2, correspondendo o limite teórico de 1 a consumo nulo fora do equipamento de TI, situação inalcançável na prática.
}

\pagina{slides/slide-46.png}{
Dois indicadores complementam o PUE em dimensões que este não capta. O WUE mede o volume de água por unidade de energia do equipamento de TI, métrica relevante em instalações que recorrem a arrefecimento evaporativo, técnica que reduz o PUE à custa de consumo hídrico e que ilustra a necessidade de avaliação multidimensional da eficiência. O CUE mede as emissões de carbono por unidade de energia do equipamento de TI, sendo determinado pela matriz energética contratada. Assinala-se que as três métricas partilham o mesmo denominador, pelo que a consolidação por virtualização, objeto da Aula 03, atua sobre todas simultaneamente ao reduzir a energia de TI necessária a uma dada carga computacional.
}

\pagina{slides/slide-47.png}{
O caso de estudo aplica os critérios da aula ao servidor Proxmox VE utilizado nos trabalhos da unidade curricular. A instalação constitui um sistema de via única em todos os subsistemas, pelo que a aplicação dos critérios do Uptime Institute a situa no nível Tier I. A execução da totalidade das máquinas virtuais da turma sobre o mesmo equipamento físico ilustra simultaneamente a partilha de recursos, analisada na Aula 01, e a existência de um ponto único de falha, matéria retomada na Aula 09. O exercício proposto consiste em estimar a disponibilidade do conjunto admitindo valores plausíveis de MTBF e MTTR para o equipamento e para a alimentação, e em identificar qual das duas grandezas condiciona o resultado obtido.
}

\pagina{slides/slide-48.png}{
A primeira metade da síntese retoma os resultados relativos à infraestrutura física e à rede. Um datacenter integra quatro subsistemas interdependentes, e o nível de redundância do conjunto é determinado pelo subsistema menos redundante, conclusão estabelecida pela questão 2. Os níveis $N$, $N+1$ e $2N$ constituem o vocabulário normalizado do qual decorre a classificação em Tiers. A topologia leaf-spine substitui a organização hierárquica em instalações com tráfego este-oeste predominante, por uniformizar a latência entre pontos internos e admitir agregação de percursos por ECMP. A razão de sobre-subscrição quantifica a contenção admitida no dimensionamento das ligações ascendentes.
}

\pagina{slides/slide-49.png}{
A segunda metade da síntese retoma os resultados quantitativos e contratuais. A disponibilidade exprime-se em função do MTBF e do MTTR, constituindo a redução do segundo, em geral, a via de melhoria com melhor relação custo-benefício, conforme quantificado na questão 5. A composição em série não excede a disponibilidade do componente menos disponível, ao passo que a composição em paralelo excede a do mais disponível, sob a hipótese de independência das falhas. Um acordo de nível de serviço é exigível na medida do seu clausulado, determinando as exclusões e o regime de compensação o seu valor efetivo. Por último, a garantia de um operador de cloud constitui um limite superior da disponibilidade observável, cuja superação exige redundância por zonas independentes.
}

\pagina{slides/slide-50.png}{
Conclui-se a exposição da Aula 02. As sete questões distribuídas pelo documento destinam-se à verificação autónoma da compreensão, e não à avaliação. A matéria aqui sistematizada não esgota a discussão conduzida em aula, constituindo antes o suporte de estudo para a aula seguinte, que abandona a infraestrutura física e passa à camada de virtualização, mecanismo pelo qual o mesmo equipamento se torna partilhável ao nível do sistema operativo.
}

\end{document}
